\documentclass[10pt,a4paper]{amsart}
\usepackage{amsmath,amssymb,amsthm,graphicx}
% For dealing with references we use the comment environment
\usepackage{verbatim}
\newenvironment{reference}{\comment}{\endcomment}
%\newenvironment{reference}{}{}
\newenvironment{slogan}{\comment}{\endcomment}
\newenvironment{history}{\comment}{\endcomment}
% For commutative diagrams we use Xy-pic
\usepackage[all]{xy}

% We use 2cell for 2-commutative diagrams.
\xyoption{2cell}
\UseAllTwocells
\theoremstyle{plain}
\newtheorem{theorem}[subsection]{Theorem}
\newtheorem{proposition}[subsection]{Proposition}
\newtheorem{lemma}[subsection]{Lemma}

\theoremstyle{definition}
\newtheorem{definition}[subsection]{Definition}
\newtheorem{example}[subsection]{Example}
\newtheorem{exercise}[subsection]{Exercise}
\newtheorem{situation}[subsection]{Situation}

\theoremstyle{remark}
\newtheorem{remark}[subsection]{Remark}
\newtheorem{remarks}[subsection]{Remarks}

\numberwithin{equation}{subsection}
% Macros
%
\def\lim{\mathop{\mathrm{lim}}\nolimits}
\def\colim{\mathop{\mathrm{colim}}\nolimits}
\def\Spec{\mathop{\mathrm{Spec}}}
\def\Hom{\mathop{\mathrm{Hom}}\nolimits}
\def\Ext{\mathop{\mathrm{Ext}}\nolimits}
\def\SheafHom{\mathop{\mathcal{H}\!\mathit{om}}\nolimits}
\def\SheafExt{\mathop{\mathcal{E}\!\mathit{xt}}\nolimits}
\def\Sch{\mathit{Sch}}
\def\Mor{\mathop{\mathrm{Mor}}\nolimits}
\def\Ob{\mathop{\mathrm{Ob}}\nolimits}
\def\Sh{\mathop{\mathit{Sh}}\nolimits}
\def\NL{\mathop{N\!L}\nolimits}
\def\CH{\mathop{\mathrm{CH}}\nolimits}
\def\proetale{{pro\text{-}\acute{e}tale}}
\def\etale{{\acute{e}tale}}
\def\QCoh{\mathit{QCoh}}
\def\Ker{\mathop{\mathrm{Ker}}}
\def\Im{\mathop{\mathrm{Im}}}
\def\Coker{\mathop{\mathrm{Coker}}}
\def\Coim{\mathop{\mathrm{Coim}}}

% Boxtimes
%
\DeclareMathSymbol{\boxtimes}{\mathbin}{AMSa}{"02}

%
% Macros for moduli stacks/spaces
%
\def\QCohstack{\mathcal{QC}\!\mathit{oh}}
\def\Cohstack{\mathcal{C}\!\mathit{oh}}
\def\Spacesstack{\mathcal{S}\!\mathit{paces}}
\def\Quotfunctor{\mathrm{Quot}}
\def\Hilbfunctor{\mathrm{Hilb}}
\def\Curvesstack{\mathcal{C}\!\mathit{urves}}
\def\Polarizedstack{\mathcal{P}\!\mathit{olarized}}
\def\Complexesstack{\mathcal{C}\!\mathit{omplexes}}
% \Pic is the operator that assigns to X its picard group, usage \Pic(X)
% \Picardstack_{X/B} denotes the Picard stack of X over B
% \Picardfunctor_{X/B} denotes the Picard functor of X over B
\def\Pic{\mathop{\mathrm{Pic}}\nolimits}
\def\Picardstack{\mathcal{P}\!\mathit{ic}}
\def\Picardfunctor{\mathrm{Pic}}
\def\Deformationcategory{\mathcal{D}\!\mathit{ef}}
\title{Lemma rieffel in brauer.tex}
\author{The Stacks Project Authors}
\date{stacks-project:a04446e57ec1fbc252a871afcec7752fb2807b14; GNU Free Documentation License 1.2 (as published in the Stacks Project COPYING)}
\begin{document}
\maketitle
\noindent\textit{Editorial note.} All mathematical statement and proof text is reproduced verbatim from the author's chapter brauer.tex of the Stacks Project, fixed at commit a04446e57ec1fbc252a871afcec7752fb2807b14. The pinned preamble is supplied as its own file; the wrapper class, the xr-hyper wiring and the \textbackslash{}externaldocument list are omitted because the assembler supplies the class.\par\medskip
\noindent\textit{Editorial note.} Numbering inside this standalone excerpt is editorial, not the full Stacks book print numbering.\par\medskip
\noindent\textit{Editorial note.} The candidate is from the corrected band-I follow-on screen (line C) of the extractor's stacks-candidates line: its author label is absent from the live master in every form searched and it belongs to none of the earlier judged pools. The H2 judgement and its reason were made in this module by reading the author's own statement and complete proof in the fixed checkout. Original candidate key: stacks:brauer.tex:lemma-rieffel.\par\medskip
\noindent\textit{Editorial note.} Reference labels in the proof resolve to navigation pointers to the original author source; these pointers do not reproduce proofs of the prerequisites.\par\medskip
\noindent\textit{Allowed standard prerequisites (editorial).}\begin{itemize}
\item Prerequisites named by the author inside the counted span. The counted span is reproduced verbatim, including the author's own citations. Its prerequisites are the earlier results of the same fixed Stacks Project version that the author cites there and in the chapter's lead-in.
\end{itemize}\medskip
\begin{lemma}
\label{lemma-rieffel}
Let $A$ be a possibly noncommutative ring with $1$ which contains no
nontrivial two-sided ideal. Let $M$ be a nonzero right ideal in $A$,
and view $M$ as a right $A$-module. Then $A$ coincides with the
bicommutant of $M$.
\end{lemma}

\begin{proof}
Let $A' = \text{End}_A(M)$, so $M$ is a left $A'$-module.
Set $A'' = \text{End}_{A'}(M)$ (the bicommutant of $M$).
We view $M$ as a right $A''$-module\footnote{This
means that given $a'' \in A''$ and $m \in M$ we have a product
$m a'' \in M$. In particular, the multiplication in $A''$
is the opposite of what you'd get if you wrote elements of $A''$
as endomorphisms acting on the left.}.
Let $R : A \to A''$ be the natural homomorphism such that
$mR(a) = ma$. Then $R$ is injective, since $R(1) = \text{id}_M$
and $A$ contains no nontrivial two-sided ideal. We claim that $R(M)$
is a right ideal in $A''$. Namely, $R(m)a'' = R(ma'')$ for $a'' \in A''$
and $m$ in $M$, because {\it left} multiplication of $M$ by any element $n$
of $M$ represents an element of $A'$, and so
$(nm)a'' = n(ma'')$ for all $n$ in $M$.
Finally, the product ideal $AM$ is a two-sided ideal, and so
$A = AM$. Thus $R(A) = R(A)R(M)$, so that $R(A)$ is a right ideal in $A''$.
But $R(A)$ contains the identity element of $A''$, and so $R(A) = A''$.
\end{proof}

\end{document}
